\begin{tabular}{@{}>{\centering\arraybackslash}p{1.2cm}>{\raggedright\arraybackslash}p{\dimexpr\textwidth-1.2cm-2.05cm-0.3cm-4\tabcolsep-2pt\relax}>{\raggedleft\arraybackslash}p{2.05cm}@{\,}>{\raggedleft\arraybackslash}p{0.3cm}@{}}
\toprule
\textbf{Round} & \textbf{Accepted change} & \multicolumn{2}{r@{}}{\textbf{Performance}} \\
\midrule
0 & \emph{Seed concept (\texttt{clockchain}), before any revision.} & 47\% & \\
1 & Propose a plan to target a specific, narrow initial user segment (e.g., digital illustrators on a single platform) for the beta launch to focus marketing and development. & 62\% & $\uparrow$ \\
2 & Propose a small-scale, planned test with a handful of digital illustrators from the target platform to gather initial workflow feedback before the beta, framing it as a step to refine the tailored interface. & 68\% & $\uparrow$ \\
3 & Propose a planned, small-scale public demonstration of the technology's verification capability using a sample file to be conducted during the crowdfunding campaign to build trust. & 80\% & $\uparrow$ \\
4 & Add a proposed, simple pricing structure (e.g., `We are considering a freemium model with X free verifications per month and a subscription for Y verifications') to the Revenue Model and Sustainability section to address a key backer question. & 82\% & $\uparrow$ \\
5 & In the `Revenue Model and Sustainability' section, propose a specific, hypothetical pricing tier (e.g., `5 free verifications/month, then \$5/month for unlimited') to make the freemium model more tangible for backers, while still noting it is subject to change based on the planned test. & 80\% & $\downarrow$ \\
6 & Add a brief, hypothetical example of a digital illustrator using the tool to prove authorship of a piece of art posted on the named social platform, to make the value proposition more tangible. & 82\% & $\uparrow$ \\
7 & Include a mock-up or wireframe screenshot of the proposed beta interface tailored for digital illustrators in the project description to visualize the solution. & 78\% & $\downarrow$ \\
8 & In the `Roadmap \& Delivery Plan' section, specify the proposed timing for the small-scale test with digital illustrators (e.g., `planned for the month following successful funding') and the public demonstration (e.g., `planned for the third week of the campaign') to add temporal clarity and show proactive planning. & 87\% & $\uparrow$ \\
9 & Define a clear, simple metric for the `small-scale, planned test' (e.g., `5 test users') to show structured planning and manageability. & 82\% & $\downarrow$ \\
10 & Define a clear, simple metric for the success of the proposed small-scale test with 5 digital illustrators (e.g., `percentage of test users who successfully completed the fingerprinting workflow without assistance') to show how feedback will be measured. & 87\% & $\uparrow$ \\
11 & Specify the exact public ledger technology (e.g., a specific blockchain, or a public timestamping service like OpenTimestamps) to be used in the technical details to enhance credibility and transparency. & 87\% & $=$ \\
12 & Detail the exact, simple metric for the small-scale test (e.g., `target: 80\% of 5 test users complete the workflow unassisted') in the Roadmap. & 88\% & $\uparrow$ \\
\bottomrule
\end{tabular}
